% !PRINTED 638
\weberpage{638}
%<<CONT
understands by $\vartheta$ a product of powers of the numbers $\varrho$ which
are derived according to the formulae (24) from the prime factors $\mathfrak{p}$
contained in $\varphi$.

\begin{center}
§. 175.

\gesperrt{The Roots of Unity in the Field $\Omega_m$.}
\end{center}

For the applications to be made further on, a more precise
knowledge of the units of the field $\Omega_m$ is required, which indeed
is also of great interest in itself. We concern ourselves
here not with the functional but only with the numerical
units, by which, as we recall, \gesperrt{are to be understood integers of the
field $\Omega_m$ whose norm is $\pm 1$,
or an integer whose reciprocal value is likewise integral.}

Among the units are certainly all the roots of unity contained
in $\Omega_m$; but we can easily prove that these
roots of unity can only be powers of $r$, with positive
and negative sign.

Namely, let $\varrho$ be a root of unity contained in $\Omega_m$, and
\weq{1}{\varrho = \varphi(r).}

If $q^h$ is the highest power of $q$ dividing the degree of $\varrho$,
then $\varrho^{q^h}$ is a root of unity whose degree is not divisible by $q$,
and
\weq{2}{\varrho^{q^h} = [\varphi(r)]^{q^h}.}

On account of the irreducibility of the cyclotomic equation in
the wider sense\footnote{Cf. Bd.~I, §.\,134 and the supplement to this volume.} all the substitutions
 $(r, r^n)$ can therefore be made in (2) without the left side
changing, and consequently $\varrho^{q^h}$ is a rational number, and since it
is a unit, it must be $= \pm 1$. Hence $\varrho$, apart from sign,
is equal to a root of unity of degree $q^h$ or, when $q = 2$,
of degree $q^{h+1}$. It is to be shown that $q^h$ or
$q^{h+1}$ cannot be greater than $m$. But this results
from (1). For if $\varrho$ were a root of unity of degree higher than $m$,
then $r$ would be a power of $\varrho$, and $\varrho$ would have to satisfy an irreducible
 equation of higher degree than $r$, and hence also of higher
degree than $\varphi(r)$, which by (1) is a contradiction. Thus
we have the theorem:
%CONT>>
